\chapter{Classical Mechanics}
\label{ch:cm}
{
\setlength\epigraphwidth{0.6\textwidth}
\epigraph {
    Motion\\
There is no motion, said a wise old Greek.\\
The other sage just smiled and walked around.\\
A better answer hardly could be found,\\
And everybody praised this tongue in cheek.\\
But, gentlemen, this funny story might\\
Remind us of a man of great renown:\\
Day after day the Sun walks up and down,\\
And yet the stubborn Galilei was right.
} {A.S. Pushkin, Translated by D. Broytman }
}

Classical mechanics  is the most important model in physics. Until the end of
the 19th century, mechanics was considered as direct and the only the basis of
all physics. By words understand or to explain a physical phenomenon meant
constructing a mechanical model of it.

Situation has changed,  other physical models exist, e.g. quantum mechanics so 
"usual" mechanics has changed its name to "classical mechanics" and still is
the central part of theoretical physics. All main physical concepts emerged from
it.

Classical mechanics is based on three principles:
\begin{description}
  \item[ Absolute Time] 
    Time is universal and global in classical mechanics. State of the system is
    determined by positions and velocities at fixed time. Every particle
    has its own coordinates however time is common. We can always measure
    time  exactly and instantly without any influence on mechanical
    systems. Time is continuous, discrete time would lead to finite
    difference equations instead of differential ones. So dynamics would be
    undefined, since one always can add periodical function to solution.
\item [Instant Interaction] All interactions are instant i.e. depend on system
    state at one fixed moment of time. It can be seen as action on a distance
        or infinite velocity of signal propagation. In particular it means that
        all accelerations and hence evolution is determined by one state (not
        necessary initial one). It is possible to think that any interaction
        takes small but finite period of time. This supposition will lead to
        functional equations instead of differential ones.
  \item [ Absolute Measurability]
    We can find out and set up exact value of any variable at any time
    instantly without disturbing the system in any way. In particular we
    can set exactly positions and velocities which means among other things
    that space is continuous.
\end{description}

Classical mechanics occupies a peculiar position among physical theories. It is,
first, a complete and self-contained subject: a body of concepts, equations, and
results describing the motion of bodies under forces, valid and useful on its own
terms, with no reference whatsoever to the Planck constant $\hbar$ or the speed of
light $c$. Celestial mechanics, rigid-body dynamics, the theory of small
oscillations, continuum mechanics, and the modern theory of dynamical systems and
chaos are all classical mechanics pursued for its own sake, and none of them is
waiting to be corrected by a more fundamental theory in order to be meaningful.

Second,  classical mechanics is
the \emph{base} on which relativistic and quantum mechanics are built. Special and
general relativity are commonly introduced as corrections to Newtonian mechanics
that become significant when velocities approach $c$ or fields become strong;
quantum mechanics is commonly introduced as a modification that becomes significant
at the scale of $\hbar$. In both cases, the classical theory is treated as a known,
well-understood limit --- the solid ground from which the generalization departs
and to which it must reduce.

This book is devoted to the second role of classical mechanics, some
interesting problems beyond this role are briefly discussed in Appendices.

Classical mechanics is usually presented in one of two logical orders:
\begin{description} 
    \item[Newton first.] The postulates are equations of motion (together with
        a catalogue of force laws and a statement about constraint reactions).
        Lagrange's equations, the principle of least action and the Hamiltonian
        formalism are \emph{derived}. This is the order of Goldstein
        \cite{GoldsteinPooleSafko2000}, Arnold's \cite{Arnold1989} first
        chapters, and most introductory courses. The strength of this approach
        is that    every step has a physical referent. Students see where the
        action comes from and, crucially, \emph{when it does not exist}.
        Constraint reactions are computed, not suppressed. The nonholonomic and
        dissipative cases are included without a change of philosophy. 
    \item[Action first.] The postulate is Hamilton's principle, $\delta S=0$ with
    $S=\int L\,dt$. Newton's equations are \emph{derived}, and conservation
        laws follow from symmetries of $L$ via Noether's theorem. This is the
        order of Landau and Lifshitz \cite{LandauLifshitz-1}, Lanczos
        \cite{Lanczos1970} and many others. Symmetry becomes the organizing
        principle: Noether's theorem explains energy, momentum and angular
        momentum at once. Coordinate independence is manifest: $\delta S=0$
        does not know what coordinates are. The road to field theory and
        quantization is short. However the phrase ``least action'' invites
        teleology. The student learns that dissipative systems are ``not
        mechanics'' and that constraint reactions are an afterthought. Most
        damagingly, the nonholonomic case produces a \emph{wrong answer} that
        is indistinguishable, in form, from the right one.
\end{description} 

Both orders lead to the same equations for simple systems of point particles,
so it is tempting to regard the choice as a matter of taste. It is not. The two
roads are \emph{not} mutually inverse once we leave the simplest class of
systems: the road from Newton to the action requires hypotheses (potential
forces, ideal constraints) that the road back silently assumes. 

In the language of chapter \ref{ch:rm}, the Newtonian laws are statements anchored
in the real world: forces are measured, constraints are realized by hinges,
rails and blades. The action is a construction of the ideal world: a functional
whose extremals are, by design, trajectories. Both routes have their strengths
and weaknesses.

Here I follow \cite{Arnold1989} as more physical, later
I shall use mainly second (action first) approach with some reservations.
I see Hamilton's principle, $\delta S=0$ not as law of nature but as 
powerful heuristic.

\section{Equations of Motion}
\label{sec:cm-eom}
In Ideal world of classical mechanics space is three-dimensional  and euclidean
$\mathbb{R}^3$ and time is one-dimensional $\mathbb{T}$.
So universe of classical mechanics is $\mathbb{T} \times \mathbb{R}^3$. 

The simplest object of classical mechanics is point mass (particle) i.e. object
whose shape, sizes and internal structure are inessential, so mass (positive
scalar $m$) is the only parameter. 
The state of a particle is fully determined by position ($\bm x \in\mathbb{R}^3$) 
and velocity ($\bm v \in\mathbb{R}^3$). 
More accurately: A motion in $\mathbb{R}^3$ is a differentiable mapping 
\[\bm x: \mathbb{I}\rightarrow \mathbb{R}^3, \mathbb{I} \subset \mathbb{R}.\] 
The derivative 
\[\dot {\bm x}(t)=\frac{d{\bm x}}{dt}(t)\]
is called the velocity vector. 
The derivative 
\[\ddot {\bm x}(t)=\frac{ d^2 {\bm x}}{dt^2}(t)\]
is called the acceleration vector. 
The image of motion (curve in $\mathbb{R}^3$) called trajectory and  curve in 
$\mathbb{T} \times \mathbb{R}^3$ which appears  as the graph of a motion, is
called a world line.

The initial state of a mechanical system (the totality of positions and
velocities of its points at some moment of time) uniquely determines all of
its motion. In particular, the initial positions and velocities determine the
accelerations. This is called `Newton's principle of determinacy'.

Let us consider the mechanical system consisting of $n$ particles 
with masses $m_j,j=1\ldots n$ and word lines $\bm x_j(t)$.
Suppose that $n$ vector functions   
\[{\bm F}_j:\mathbb{R}^{3 n}\times\mathbb{R}^{3 n}
\times\mathbb{R}\rightarrow\mathbb{R}^3,j=1\ldots n\]
are defined. Then the motion of all particles is determined by equations
\begin{equation}
  \label {newton-eq}
    m_j \ddot {\bm x}_j={\bm F}_j({\bm x}_k,\dot {\bm x}_k,t) \quad j,k=1\ldots n 
\end{equation}
and initial conditions
\begin{equation}
\label {newton-ic}
    {\bm x}_j(t_0)={\bm y}_j\quad\dot {\bm x}_j(t_0)={\bm v}_j  \quad j=1\ldots n 
\end{equation}

Newton's equation \eqref{newton-eq} is  the basis of mechanics. Functions 
$\bm F_j$ are called forces they should be determined  for each specific 
mechanical system. If our system is isolated i.e. includes explicitly all 
interacting particles, then galilean invariance puts some restrictions 
on forces.

Namely 
\begin{equation*}
     m_j \ddot {\bm x}_j=
    {\bm F}_j({\bm x}_j-{\bm x}_k,\dot {\bm x}_j-\dot {\bm x}_k)\quad j,k=1\ldots n 
\end{equation*}
practically such systems occur in celestial mechanics. In most other cases the
system considered is not isolated and forces of general kind \eqref{newton-eq} 
appear as a result of interaction with environment.

Some forces can not be set explicitly they appear as a result of constraints
imposed on the system and should be figured out form the prescribed movement.
\begin{equation}
  \label {newton-eq-const}
    m_j \ddot {\bm x}_j={\bm F}_j({\bm x}_k,\dot {\bm x}_k,t) 
    +{\bm R}_j \quad j,k=1\ldots n 
\end{equation}
These constraint forces ${\bm R}_j$ are result of collective action of 
\emph{all constraints}. If a system consists of mass points connected by 
rods, hinges, etc., then the need may arise to talk about the constraint force 
of some particular constraint. However the concept of a constraint force for an 
individual constraint is generally \emph{impossible to define}.

Mathematics of  Cauchy problem (\ref{newton-eq},\ref{newton-ic}) 
one can find in e.g. \cite{Arnold1992}. Generally the unique solution exists if
all functions are smooth enough. However smoothness may be violated in 
physically important cases, anyway functions $\bm x_j(t)$ should be continuous,
$\bm v_j(t)$  may have finite jumps as a results of strong force acting for 
very short time (e.g. elastic collision).

Note that initial conditions have nothing to do with causality. Causality makes
sense in Real world only, in Ideal world we are free to set final conditions or
fix $\bm x_j(t_1)$ and $\bm v_j(t_1)$ for some intermediate time $t_1$. In
many cases we are interested in solution of equation \eqref{newton-eq} with
fixed initial velocities and final coordinates  (aiming) or fixed coordinates
on both ends (boundary problem). Anyway when conditions are set at two
different times  solution may be not unique (unless force is changing
relatively slow) \footnote {I discovered this fact around 1980 and was very
proud of myself until I found papers \cite{Maslov1962a} and \cite{Maslov1962b}.}.

If we allow only potential forces i.e. 
\begin{equation}
  \label {pot-force}
        \bm F_j(\bm x_k,t)=-\frac{\partial U(\bm x_k,t)}{\partial \bm x_j} \quad j,k=1\ldots n 
\end{equation}
we shall see that equations of motion (\ref{newton-eq}) can be written as 
\begin{equation}
   \label {lagr-eq}
    \frac{d}{dt} \frac{\partial L({\bm x}_k,{\bm v}_k,t)}{\partial {\bm v}_j}= 
    \frac{\partial L({\bm x}_k,{\bm v}_k,t)}{\partial {\bm x}_j}\quad j,k=1\ldots n 
\end{equation} 
where 
\begin{equation}
  \label{lagr-nat}
        L({\bm x}_k,{\bm v}_k,t)=\Sigma_1^n \frac{m_k}{2} {\bm v}_k \cdot {\bm v}_k-U({\bm x}_k,t) 
\end{equation}
Function $L({\bm x}_k,{\bm v}_k,t)$ is called Lagrangian, vectors 
\[{\bm p}_j= \frac{\partial L({\bm x}_k,{\bm v}_k,t)}{\partial {\bm v}_j}
=m_j \bm v_j \] 
are called momenta. 
We can write the equations \eqref{lagr-eq} in more symmetrical way:
\begin{equation}
   \label {ham-eq}
   \begin{split}
       \dot {\bm x}_j&=\frac{\partial H(\bm x_k,\bm p_k,t)}{\partial \bm p_j}  \\ 
        \dot {\bm p}_j&= -\frac{\partial H(\bm x_k,\bm p_k,t)}{\partial \bm x_j}\quad j,k=1\ldots n    
   \end{split}
\end{equation}
where 
\begin{equation}
        H(\bm x_k,\bm p_k,t)=\Sigma_1^n \frac{1}{2 m_k} \bm p_k \cdot \bm p_k+U(\bm x_k,t) \label{ham-nat}
\end{equation}
Note that functions $L$ (Lagrangian) and $H$ (Hamiltonian) are connected by Legendre transform:
\begin{equation}
        L(\bm x_k,\bm v_k,t)=\Sigma_1^n  \bm p_k \cdot \bm v_k -H(\bm x_k,\bm p_k,t)|_{\bm p_k=m_k\bm v_k}
\quad k=1\ldots n \label {lagr-ham}
\end{equation}


\begin{definition}
        The equation
        \[\frac{d}{d t}(\frac{\partial L}{\partial \dot {\bm x}})-\frac{\partial L}{\partial  {\bm x}}=0 \]
        is called the Euler-Lagrange equation for the functional
        \[S=\int_{t_1}^{t_2} L(\bm x,\dot {\bm x},t) dt \]
\end{definition}
Let \( \gamma= \bm x(t) \quad t_1 \le t \le t_2 \) be a curve  in  $\mathbb{R} \times \mathbb{R}^n$ and 
$L: \mathbb{R}^n \times \mathbb{R}^n \times \mathbb{R}\rightarrow \mathbb{R}$ a function of $2 n +1$
variables, then
\begin{theorem}
    The curve $\gamma$ is an extremal of the functional \[S=\int_{t_1}^{t_2} L(\bm x,\dot {\bm x},t) dt \]
        on the space of curves joining $(t_1, \bm x_1)$ and $(t_2,\bm x_2)$, if and only if the Euler-
Lagrange equation is satisfied along $\gamma$.
\end{theorem}
This is a system of n second-order equations, and the solution depends on
$2 n$ arbitrary constants. The $2 n$ conditions $\bm x(t_1) = \bm x_0$ and $\bm x(t_2) = \bm x_2$ are used
for finding them.
Two important remarks:
\begin{enumerate}
        \item There may be many extremals connecting two
given points or none at all.
\item The condition for a curve  to be an extremal of a functional does not depend
on the choice of coordinate system. This paves way to geometrisation of mechanics.
\end{enumerate}

It is easy to see that equation of motion (\ref{lagr-eq}) are actually
Euler-Lagrange equations for lagrangian \eqref{lagr-nat}. So we have  Hamilton's
principle of least action.

Equations of motion (EOM) are the heart of classical mechanics in both its roles.
However equations themselves are not enough, to get meaningful solution
one should fix the \emph{problem} i.e. specify the conditions imposed.
The problems usually  associated with classical mechanics are:
\begin{description}
    \item [IVP / Cauchy problem]  The full state (position and velocity, or a
  phase-space point) at one instant $t_0$ is specified.
   The Picard--Lindel\"of theorem gives
  local existence and uniqueness for Lipschitz forces; global existence
  requires a no-blow-up bound.
   This is what time-stepping integrators  actually solve.
    Two large class of methods are in use: 
    \begin{enumerate} 
    \item "Usual" one-time integrators when differential equations are replaced
        with finite difference ones defining snapshots at discrete time steps.
        These integrators allow accurate description of standard problem:
        initial conditions, finite time interval. 
    \item Geometric integrators which are less accurate for any single trajectory
    but accurately present qualitative picture for large set of initial data
        and long times. Such integrators may calculate  coordinates and momenta
        at different times
    \end{enumerate}
    One can look for details in \cite{HairerLubichWanner2006} for example.

    \item [Boundary Value Problem (BVP)] Data at two or more distinct times/points,
  e.g. $x(t_0) = x_0$, $x(t_1) = x_1$ --- not a full state at one
  instant is fixed.
   Existence and uniqueness are \emph{not} automatic (unlike the
  IVP): zero, one, or many solutions can occur.
   \textbf{Prototypes:} the brachistochrone problem, geodesics
  between two points, Lambert's problem, and any Euler--Lagrange
  equation arising from a fixed-endpoint action principle.
   \textbf{Numerically:} shooting methods, relaxation /
  finite-difference methods, collocation.

    \item [Scattering Problem] An asymptotic free state as $t \to
  -\infty$ (or $r \to \infty$) is specified.
   The asymptotic outgoing state as $t \to
  +\infty$ --- the $S$-matrix or cross-section --- given a localized
  interaction region in between is sought.
   This is asymptotic BVP: the conditions sit at $t = \pm\infty$ rather
   than at finite times, which   is exactly why it needs its own machinery 
   ($S$-matrix, asymptotic
  completeness) rather than shooting/relaxation methods.
 Often the trajectory itself is discarded --- only the
  in $\to$ out map matters.

    \item [Inverse Problem] Having observed trajectories or scattering data,
   one searches  the force law / potential / EOM itself
  (inverse scattering; Newton inferring $1/r^2$ from Kepler's laws is
  the archetype).
\end{description}

\section{Covariance, Invariance and Symmetry}
\label{sec:cm-cis}
\epigraph{‘When I use a word,’ Humpty Dumpty said in rather a scornful tone,
‘it means just what I choose it to mean --- neither more nor less.’} {Lewis
Carroll}

Few words in physics are used as loosely as \emph{invariant}, \emph{covariant}
and \emph{symmetric}. One reads of the ``Lorentz invariance of Maxwell's
equations'' and, a page later, of their ``Lorentz covariance''; of ``general
covariance'' as the foundation of general relativity and, in the next
paragraph, as a triviality of notation; of ``gauge symmetry'' that is plainly
not a symmetry in the sense in which rotations are. The three notions are
tied together by one logical structure, but they live at \emph{different
levels of description}, and mixing the levels is a dependable source of
confusion.

\begin{itemize}
\item \textbf{Invariance} is a property of an \emph{object}: a number, a
  function, a tensor field, a functional, a particular solution. The object is
  left unchanged by a transformation.
\item \textbf{Covariance} is a property of a \emph{formulation}: a family of
  quantities, or an equation built from them, transforms in a definite,
  prescribed way (a representation of the transformation group), so that the
  equation keeps its form.
\item \textbf{Symmetry} is a property of a \emph{law} (a dynamical system): a
  transformation that carries solutions into solutions, or leaves the action
  unchanged.
\end{itemize}

Symmetry is the only one of the three that is a claim about the real world. It
says that if an experiment on an isolated system is repeated after the
\emph{entire} arrangement has been displaced, rotated or boosted, the outcome
is the transformed outcome. That claim can be, and has been, tested with
extraordinary precision (Michelson--Morley-type experiments for boosts,
Hughes--Drever-type experiments for rotations). Covariance, by contrast, is
a statement about how we write equations; whether it carries any physical
content depends on what else is held fixed.

A \emph{passive} transformation relabels the points of spacetime, $x\mapsto
x'=f(x)$, and asks how the \emph{same} physical quantity is described in the
new coordinates. An \emph{active} transformation by $g$ moves the quantity
itself, keeping the coordinates.The new quantity \emph{function} is given by
the same formula in both cases what differs is the interpretation. In the
passive reading one compares two descriptions of one physical situation; in the
active reading one compares two different physical situations in one
description. Symmetry is an active notion: it asks whether the second situation
is also allowed by the law. Covariance is naturally passive: it asks whether
the \emph{description} keeps its form. The active reading carries a warning
that will turn out to be the heart of the matter: \emph{everything that enters
the equations must be transformed, or else be invariant.}

\begin{example}[Coulomb field of a fixed charge]
    \label{ex:coulomb} The equation
    $\nabla\!\cdot{\bm{E}}=4\pi\rho$ with a prescribed source
    $\rho=q\,\delta^3({\bm{x}}-{\bm{x}}_0)$ has the solution
    ${\bm{E}}({\bm{x}})=q({\bm{x}}-{\bm{x}}_0)/|{\bm{x}}-{\bm{x}}_0|^3$.
    Translating the field alone, ${\bm{E}}_a({\bm{x}})={\bm{E}}({\bm{x}}-\bm
    a)$, gives a field whose divergence is $4\pi
    q\,\delta^3({\bm{x}}-{\bm{x}}_0-\bm a)$: it is \emph{not} a solution for
    the original source. If the source is translated too, the pair
    $({\bm{E}}_a,\rho_a)$ is a solution. Translation is therefore a symmetry of
    the closed system ``field $+$ charge,'' and is broken as soon as the charge
    is treated as a fixed, external, non-dynamical object. 
\end{example}

Whatever is held fixed and not transformed is called a \emph{background} (or
\emph{absolute}) structure: a prescribed source, a fixed metric, a preferred
velocity field. Prescribing backgrounds is unavoidable in practice, and it is
precisely what separates covariance from symmetry.


\paragraph{Covariance and invariance of Newton's law.}
Take $N$ particles with $m_a\ddot{{\bm{x}}}_a=\sum_{b\ne a}\bm F_{ab}$, where
$\bm F_{ab}=-\nabla_{a}V(|{\bm{x}}_a-{\bm{x}}_b|)$ is a central force. The Galilei group
(10 parameters) acts as
\begin{equation}\label{galilei}
  t'=t+s,\qquad {\bm{x}}'=R\,{\bm{x}}+{\bm{v}}\,t+\bm a,\qquad R\in SO(3).
\end{equation}
For \eqref{galilei},
\begin{equation}
  \ddot{{\bm{x}}}'_a=R\,\ddot{{\bm{x}}}_a,\qquad
  {\bm{x}}'_a-{\bm{x}}'_b=R\,({\bm{x}}_a-{\bm{x}}_b).
\end{equation}
The acceleration is thus \emph{invariant} under boosts and translations and
\emph{covariant} under rotations, and boosts and translations cancel in
differences of positions. Hence $\bm F'_{ab}=R\bm F_{ab}$ and the equation of
motion is covariant. The velocity $\dot{\bm{x}}'=R\dot{\bm{x}}+{\bm{v}}$, by contrast, is
\emph{not} covariant under boosts (it transforms inhomogeneously), so it can
enter a boost-symmetric force law only through differences of velocities.

\paragraph{Breaking by a background: the medium.}
Add friction against a medium with prescribed velocity field $\bm w({\bm{x}},t)$:
\begin{equation}\label{friction}
  m\ddot{\bm{x}}=-\gamma\,(\dot{\bm{x}}-\bm w).
\end{equation}
Under a boost, $\bm w\to\bm w'=\bm w+{\bm{v}}$, so \eqref{friction} is
\emph{covariant} under the Galilei group provided $\bm w$ is transformed
along with ${\bm{x}}$. But the boost is a symmetry only if the background is
unchanged. For a medium at rest ($\bm w=0$) the stabilizer consists of time
translations, space translations and rotations: boosts are excluded. The medium
plays the role of Anderson's absolute object, and its rest frame is a
preferred frame. Including the medium's atoms as dynamical degrees of
freedom restores the boost symmetry of the \emph{closed} system, exactly as in
Example~\ref{ex:coulomb}.

\paragraph{Boosts: invariance up to a total derivative.}
The Lagrangian $L=\sum\tfrac12m_a\dot{\bm{x}}_a^2-V$ is not invariant under boosts
but is \emph{quasi-invariant}: for the infinitesimal boost
$\delta{\bm{x}}_a=\eps\,\bm u\,t$,
\begin{equation}
  \delta L=\eps\sum_a m_a\bm u\cdot\dot{\bm{x}}_a=\eps\,\frac{\dd}{\dd t}\Big(\bm u\cdot\sum_a m_a{\bm{x}}_a\Big),
\end{equation}
which has the form \eqref{quasi} with $K=\bm u\cdot\sum m_a{\bm{x}}_a$. The
Noether charge $Q=\sum_a{\bm{p}}_a\!\cdot\delta{\bm{x}}_a/\eps-K=-\bm u\cdot\big(M\bm X-\bm P\,t\big)$
shows that
\begin{equation}
  \bm G=M\bm X_{\rm cm}-\bm P\,t=\text{const.}
\end{equation}
which says that the centre of mass moves uniformly. The quasi-invariance has a
deeper consequence: $\pb{G_i}{P_j}=M\delta_{ij}$, so the total mass $M$
appears as a \emph{central charge} of the Galilei algebra, in contrast to
the Poincar\'e algebra where $[K_i,P_j]\propto\delta_{ij}H$. This is why mass
is a conserved parameter of Galilean systems and why its quantum version
requires a projective representation.

\paragraph{Non-inertial frames and the geometric formulation}
Under a time-dependent change of frame ${\bm{x}}'=R(t){\bm{x}}+\bm a(t)$ Newton's law
acquires extra terms (Coriolis, centrifugal, Euler, and translational
inertial forces) and no longer has the same form. So the law is covariant under the Galilei group
\eqref{galilei} only. It can, however, be \emph{reformulated} so that it is
covariant under arbitrary coordinate changes: this is Cartan's
Newton--Cartan geometry. Free fall is the geodesic equation
of a connection with $\Gamma^i{}_{00}=\partial_i\Phi$, and the absolute
objects are the time one-form $t_\mu=\partial_\mu t$ and the spatial
cometric $h^{\mu\nu}$. The formulation is generally covariant, the symmetry
is the subgroup preserving $t_\mu$ and $h^{\mu\nu}$, which is larger than the
Galilei group (it includes ${\bm{x}}\to{\bm{x}}+\bm a(t)$ acting on $\Phi$ by a term
linear in ${\bm{x}}$, proportional to $\ddot{\bm a}$, an expression of the
Newtonian equivalence principle). 

Symmetry is often introduced to students as a single, monolithic idea: ``the
Lagrangian is invariant under such-and-such transformation, therefore there is
a conserved quantity, therefore the system respects the symmetry.'' In
practice, the situation is considerably richer, and the richness is not a
technicality but the source of some of the deepest phenomena in physics, from
the accidental degeneracies of the Kepler problem to the existence of the
electron mass and the pion.

It is useful to keep distinct objects in view:
\begin{enumerate}
\item the \textbf{Lagrangian} (or action, or Hamiltonian) that \emph{defines} 
    the dynamics;
\item the \textbf{equations of motion} (EOM) that it \emph{generates}; and
\item an actual \textbf{solution} of those equations, singled out by initial
    or boundary data;
\item the state that is actually realized as the ground state or equilibrium 
    configuration of the system.\footnote { Spontaneous Symmetry Breaking: 
        The Lagrangian, the equations of motion,
        \emph{and} the boundary conditions are all fully symmetric, and yet the
        stable, physically realized solution is not. This is not contingent: it
        is a structural consequence of the existence of a degenerate manifold
        of energy-minimizing (or free-energy-minimizing) configurations,
        together with an instability of the symmetric configuration and, in the
        case of a genuine phase transition, the thermodynamic (infinite-volume)
        limit.}

\end{enumerate}
Each of these objects has its own symmetry group, and the three groups need not
coincide. 

Let $q^a(t)$, $a=1,\dots,n$, be generalized coordinates on a 
configuration
space $Q$, with Lagrangian $L(q,\dot q,t)$ and action
\[
S[q] = \int_{t_1}^{t_2} L(q(t),\dot q(t),t)\,dt .
\]
\begin{equation*}
   p_a=\frac{ \partial L(q,\dot q,t)}{\partial \dot q_a}
\end{equation*}
The Euler--Lagrange (EL) operator acting on $L$ is
\[
[L]_a = \frac{\partial L}{\partial q^a} - 
\frac{d}{dt}\frac{\partial L}{\partial \dot q^a}
\]
and the equations of motion are $[L]_a=0$.

\begin{definition}[Variational / Noether symmetry]
An infinitesimal transformation $\delta q^a = \epsilon X^a(q,\dot q,t)$ is a
    \emph{variational symmetry} of $L$ if there exists a function $F(q,\dot
    q,t)$ such that, \emph{identically in $q(t)$} (not only on solutions),
    $\delta L=\epsilon \tfrac{dF}{dt}$. Equivalently, the action changes only by
    a boundary term, $\delta S = \epsilon\big[F\big]_{t_1}^{t_2}$, for
    arbitrary paths, not just extremal ones.
\end{definition}

Evidently adding the total derivative
\begin{equation*}
    \frac{d F (q,t)}{d t}=\frac{\partial F(q,t)}{\partial t} + \frac{ \partial F (q,t)}{\partial q_j} \dot q_j
\end{equation*}
to Lagrangian does not change equations of motion. However canonical momenta
become
\[ P_a=p_a+\frac{\partial F}{\partial q_a} \]
it is nothing but gauge invariance. Transformation $ p,q \rightarrow P,q $ is evidently
canonical and new Hamiltonian becomes 
\[ H_{new}=H(p(P),q)- \frac{\partial F}{\partial t}.\]

We are used to think that gauge invariance is about electromagnetism, actually it appears
already in classical mechanics. 
So canonical momenta have no physical meaning (since it is not gauge invariant).
This fact correspond to geometrical point of view: $p_a,q_a$ are merely local
coordinates on phase space.

Having EOM as system of ordinary differential equations,
we can transform it into equivalent one, say by multiplying every equation by arbitrary 
function which never becomes zero. This new system may correspond to another Lagrangian and
the difference between new and old one need not come down to total derivative. Most probably
new system will not be Lagrangian at all.

Let us look at two-dimensional harmonic oscillator.
Usual Lagrangian is
\begin{equation}
   \label {ho1}
    L_1= \frac{1}{2} (\dot x^2+\dot y^2) -\frac{ \omega^2}{2} ( x^2+ y^2)
\end{equation}
which yields equation  of motions:
\begin{equation}
   \label {hoe1}
   \begin{split}
       \ddot x+ \omega^2 x & =0 \\
       \ddot y+ \omega^2 y & =0 
   \end{split}
\end{equation}
Same equations of motion come form 
another Lagrangian
\begin{equation}
   \label {ho2}
    L_2= \frac{1}{2} (\dot x^2-\dot y^2) -\frac{ \omega^2}{2} ( x^2- y^2)
\end{equation}
or
\begin{equation}
   \label {ho3}
    L_3= \dot x\dot y - \omega^2  x y .
\end{equation}
Evidently these Lagrangians are not connected by adding total derivative.
Corresponding  Hamiltonians 
\begin{equation}
   \label {ho2h}
    H_2= \frac{1}{2} ( p_x^2-p_y^2) +\frac{ \omega^2}{2} ( x^2- y^2)
\end{equation}
and
\begin{equation}
   \label {ho3h}
    H_3= p_x p_y +\omega^2  x y
\end{equation}
are not even positive definite!

\begin{definition}[Symmetry of the equations of motion]
A transformation $g$ (finite or infinitesimal) is a \emph{symmetry of the EOM},
    or a \emph{dynamical symmetry}, if it maps every solution of $[L]_a=0$ to 
    another solution of the \emph{same} equations, i.e. it maps the solution 
    set to itself. No reference to the action or to boundary terms is required 
    by this definition.
\end{definition}

Two inclusions organize the entire subject:
\begin{itemize}
    \item $G_L \subseteq G_{EOM}$: \emph{every symmetry of the Lagrangian is a
        symmetry of the equations of motion, but the equations of motion can
        admit strictly more symmetry than is manifest in the Lagrangian one
        wrote down.} 
    \item $G_{sol} \subseteq G_{EOM}$: \emph{a particular solution, or a
        particular physically realized state, generally respects only a
        subgroup of the symmetry of the equations that produced it.} 
        Asymmetry is put in by hand
        through the choice of initial or boundary data. This is contingent:
        choosing symmetric data restores a symmetric solution, and averaging
        over the group orbit of solutions restores the full symmetry at the
        ensemble level.
   \end{itemize}

\begin{example}[Free particle: hidden $SL(2,\mathbf R)$]
The Lagrangian $L=\tfrac12 m\dot x^2$ manifestly has the Galilei symmetry
    (translations, boosts) plus time translation. But the equation of motion
    $\ddot x=0$ has straight lines as its general solution, and the space of
    straight lines is left invariant by the larger, one-dimensional-time
    conformal (``Schr\"odinger/M\"obius'') group acting projectively on $t$,
\[
t \;\longmapsto\; \frac{\alpha t+\beta}{\gamma t+\delta},
    \qquad \alpha\delta-\beta\gamma=1,
\]
together with a compensating rescaling of $x$. This $SL(2,\mathbf R)$ ``dynamical''
    symmetry (closely related to the conformal/Schr\"odinger symmetry of the
    free Schr\"odinger equation, and to the Ermakov--Lewis invariant of the
    time-dependent oscillator) is not visible as a point symmetry of
    $L=\tfrac12 m\dot x^2$ written in the ordinary way, yet it is a genuine
    symmetry of $\ddot x=0$.
\end{example}

\begin{example}[Kepler problem: hidden symmetry]
For $H={\bm{p}}^2/2m-k/r$ the geometric symmetry is $SO(3)\times\mathbb R_t$
(rotations and time translation), with conserved ${\bm{L}}={\bm{x}}\times{\bm{p}}$ and $H$.
There is, however, an additional conserved vector, the Runge--Lenz vector
\begin{equation}\label{rl}
  \bm A={\bm{p}}\times{\bm{L}}-mk\,\hat{{\bm{x}}},\qquad \pb{A_i}{H}=0 .
\end{equation}
The transformation it generates, $\delta{\bm{x}}=\eps\,\pb{{\bm{x}}}{{\bm{n}}\!\cdot\!\bm A}$,
\begin{equation}
  \delta{\bm{x}}=\eps\,\big[\,2({\bm{n}}\!\cdot\!{\bm{x}})\,{\bm{p}}-({\bm{x}}\!\cdot\!{\bm{p}})\,{\bm{n}}-({\bm{n}}\!\cdot\!{\bm{p}})\,{\bm{x}}\,\big],
\end{equation}
involves the \emph{momentum}: it is not a point transformation of
configuration space and is not related to any covariance of the
equation $m\ddot{\bm{x}}=-k{\bm{x}}/r^3$ under coordinate changes. It is nonetheless a
canonical symmetry of the Hamiltonian, and the Poisson algebra
\begin{equation}
  \pb{L_i}{L_j}=\epsilon_{ijk}L_k,\quad
  \pb{L_i}{A_j}=\epsilon_{ijk}A_k,\quad
  \pb{A_i}{A_j}=-2mH\,\epsilon_{ijk}L_k
\end{equation}
closes on $so(4)$ for $H<0$ (after the rescaling $\bm A\to\bm A/\sqrt{-2mH}$),
$e(3)$ for $H=0$, and $so(3,1)$ for $H>0$. The extra symmetry explains why
bound orbits are closed (the perihelion does not precess) \cite{goldstein}
and, in the quantum problem, the extra degeneracy of the hydrogen levels.
It is an example of a \emph{dynamical} (hidden) symmetry:
a symmetry of the law that is not visible as a covariance of the formulation.
\end{example}

\begin{example}[Integrable systems: toroidal symmetry from action-angle variables]
More generally, for a Liouville-integrable Hamiltonian system with $n$ degrees
    of freedom, one can pass to action-angle variables $(I_i,\theta^i)$ in
    which $\mathcal{H} = \mathcal{H}(I)$ depends only on the actions. 
    The flow generated by
    $\mathcal{H}$ is then, trivially, an ($n$-torus of) simultaneous phase shifts
    $\theta^i \to \theta^i+c^i$, an abelian $U(1)^n$ symmetry group of the
    equations of motion. This symmetry is exact and completely general for
    integrable systems, yet it is essentially never visible as a point symmetry
    of the Lagrangian written in the original coordinates; the free-particle
    $SL(2,\mathbf R)$ and the Kepler $SO(4)$ above are, in this sense, special
    cases of the general toroidal symmetry structure of integrable systems,
    made visible only after a nontrivial (often nonlocal, momentum-dependent)
    change of variables.
\end{example}

\paragraph{The stabilizer of a solution}
Suppose $G_{EOM}$ acts on the solution space $\mathcal S$ of $[L]_a=0$.
 this action is well defined at least for
$g\in G_L$, and by definition for all of $G_{EOM}$). A single solution
$q(t)\in\mathcal S$, picked out by initial data $(q(t_0),\dot q(t_0))$ or by
boundary data, need not be a fixed point of this action. Its symmetry is the
\emph{stabilizer subgroup}
\[
    G_{sol}=\mathrm{Stab}_{G_{EOM}}(q) = \{ g\in G_{EOM} : g\cdot q
    = q \} \subseteq G_{EOM} .
\]
Generically $G_{sol} \subsetneq G_{EOM}$: the group moves the given
solution to a \emph{different} member of $\mathcal S$ (still a valid solution —
the equations do not care — but a different trajectory), rather than mapping it
to itself.

\begin{example}[Central force problem]
The Kepler/central-force EOM has full $SO(3)$ symmetry.
     A single
    orbit, however, has a definite, fixed angular momentum vector $\bm L = \bm
    r\times\bm p\neq 0$, determined by the initial position and velocity, and
    lies entirely in the plane orthogonal to $\bm L$. The stabilizer of this
    orbit inside $SO(3)$ is only the $SO(2)$ subgroup of rotations about the
    $\bm L$ axis (together with, for a closed orbit, possibly a discrete
    reflection). The other two rotational generators move the orbital plane to
    a different plane, i.e. map this solution to a genuinely different
    solution. Nothing is wrong with the dynamics; the reduction from $SO(3)$ to
    $SO(2)$ is entirely a consequence of the choice of initial data.
\end{example}

\begin{example}[Free particle]
$\ddot x=0$ is invariant under all spatial translations and Galilean boosts. A
    single trajectory $x(t)=x_0+v_0t$ is left invariant only by time
    translations (and, trivially, the identity among boosts/translations); a
    spatial translation or a boost maps it to a different straight line.
    Symmetric \emph{initial data} — e.g. $x_0=0,\,v_0=0$, the particle
    permanently at rest at the origin — is the unique choice whose stabilizer
    is the full translation-and-boost group; any other data breaks it
    explicitly.
\end{example}

\paragraph{Summary}
\begin{enumerate}
    \item \textbf{``Covariant'' and ``invariant'' are not synonyms.} The fields
  ${\bm{E}},{\bm{B}}$ are covariant (they mix under boosts); $E^2-B^2$ and
  ${\bm{E}}\!\cdot\!{\bm{B}}$ are invariant. A covariant equation $\mathcal E=0$ is
  \emph{form-invariant}, while its left-hand side is not invariant but
  transforms in a representation. The standard phrase ``the equations are
  Lorentz invariant'' is acceptable shorthand only because it is clear from
  context that it is the \emph{solution set} that is invariant.
\item \textbf{Covariance is not physics; its stabilizer is.} Any equation can
  be written covariantly; what matters is
  which structures must then be declared as backgrounds, and whether they are
  invariant.
\item \textbf{A symmetry of the equations need not be a symmetry of the action.}
  The rescaling $x\to\lambda x$ of a \emph{linear} equation, such as the
  harmonic oscillator, maps solutions to solutions, but changes the action by
  $\lambda^2$; it has no Noether charge. Electromagnetic duality is another case.
  The converse (variational $\Rightarrow$ equation symmetry) holds
  (Proposition~\ref{prop:variational}).
\item \textbf{A symmetry of the law is not a symmetry of the solution.}
  Elliptic orbits, the Coulomb field, the vacuum of a broken theory, and the
  Schwarzschild metric all have smaller symmetry than their laws.
\item \textbf{Gauge ``symmetry'' is a redundancy.} Constant gauge
  transformations are symmetries with charges; local gauge transformations are
  not, and ``spontaneous breaking of a local symmetry'' is a statement about
  the \emph{description} and not about physical degeneracy.
\item \textbf{Transform everything.} An active symmetry transformation must act on
  the whole closed system (fields \emph{and} sources, apparatus,
  boundary conditions). Leaving a source untransformed amounts to introducing a
  background that breaks the symmetry (Example~\ref{ex:coulomb}).
\end{enumerate}

\section{Finite and infinite numbers of degrees of freedom}
\label{sec:cm-dof}
Usually mechanical systems (finite number of degrees of freedom) 
and field theories (infinite number of degrees of freedom) are seen
as absolutely different. However form symmetry (or covariance) point
of view the two are
not different theories: a mechanical system is a field theory on a
one-dimensional base manifold, and a field theory is a mechanical system whose
coordinates carry a continuous label. This is an illustration of the strength
of Action First approach.

\paragraph{Mechanics as a field theory on $M=\mathbb R_t$.}
Take the base manifold to be the time axis, with the single coordinate
$x^0\equiv t$, and identify
\begin{equation}
    \label{dictionary-id}
    \begin{split}
        \varphi_a(x^\mu)&\leftrightarrow q^i(t),\\
        \mathcal L(\varphi_a,\partial_\mu\varphi_a;x)&\leftrightarrow L(q^i,\dot q^i,t),\\
        \partial_\mu &\leftrightarrow \frac{\dd}{\dd t}.
    \end{split}
\end{equation}
In $d=1$ the divergence $\partial_\mu J^\mu$ is a single derivative, so the
current \emph{is} the charge and the conservation law is constancy in time.
Transformations are reparametrizations of time and the symmetry is time 
translation. Noether's theorem gives conservation of $H$. 

\paragraph{Field theory as a mechanical system with a continuum of degrees of
freedom.}
In the opposite direction, a $3{+}1$ split of a field theory turns it into a
mechanical system: set 
\[q_{\bm{x}}(t)=\varphi(\bm{x},t),\]
so that the discrete index
$i$ becomes the continuous label $\bm{x}$. Then
\begin{equation*}
    \begin{split}
        L&=\int\dd^3 x \mathcal{L} \\
        \pi(\bm{x})&=\frac{\delta \mathcal{L}}{\delta\dot\varphi(\bm{x})}.
    \end{split}
\end{equation*}
A chain of masses is a textbook example. On a ring of $N$
masses,
\begin{equation*}
  L=\sum_{n}\Big[\tfrac12m\,\dot q_n^2-\tfrac12k\,(q_{n+1}-q_n)^2\Big].
\end{equation*}
Set $q_n(t)=\varphi(na,t)$, $m=\rho a$, $k=\tau/a$, and keep the length $Na$
fixed. As $a\to0$,
\begin{equation*}
  L\ \to\ \int\dd x\,\Big[\tfrac12\rho\,\dot\varphi^2-\tfrac12\tau\,\varphi'^2\Big],
  \qquad\rho\ddot\varphi=\tau\varphi'' .
\end{equation*}
This is an example of 'singular limit'. Sure the analogy of mechanics and
filed theory is pure formal and does not mean that physics is the same.

\section{Geometric Setup}
\label{sec:cm-geom}

The configuration space of a mechanical system comprising $k$ particles is 
\[   
 \underbrace{\mathbb R^3\times\mathbb R^3\times\cdots\times\mathbb R^3}_{k\ \text{times}}. 
\]
which is naturally isomorphic (as vector space, topological space, and smooth manifold) to
\begin{equation}
   \label {flat-Q}
   Q=\mathbb R^{3k}.
\end{equation}

However this isomorphism hides the important difference ---
the \(3k\) coordinates are not just arbitrary coordinates. They come in \(k\)
distinguished groups of three. Rotations illustrate the difference:
Suppose we rotate physical space by
\[ R\in SO(3). \]
For \(k\) particles, the same physical rotation acts on every particle:
\[
    (\mathbf r_1,\ldots,\mathbf r_k) \longrightarrow (R\mathbf r_1,\ldots,R\mathbf r_k). 
\]
Thus the physical rotation group is represented as
$ SO(3)\subset SO(3k)$, with the action
\[ R\mapsto 
\begin{pmatrix} R&0&\cdots&0\\ 
0&R&\cdots&0\\ 
\vdots&\vdots&\ddots&\vdots\\
0&0&\cdots&R 
\end{pmatrix}. 
\]
Natural rotation group of configuration space \eqref{flat-Q} is \(SO(3k)\).
But a general element of \(SO(3k)\) can mix coordinates belonging to different particles,
such a transformation is generally not a rotation of physical three-dimensional space.

This distinction becomes particularly important in configuration-space
formulations of  quantum mechanics, because it is closely related
to the question of whether the \(3k\)-dimensional configuration space should be
regarded as a merely mathematical construction or as having some independent
physical reality.
\footnote{In quantum mechanics  $ 1,2,\ldots,3k $,
are not physically meaningful for identical particles. One should then identify
configurations related by a permutation,
so the physical configuration space is more properly
\[ Q_{\rm identical} = (\mathbb R^3)^k/S_k. \]}

Flat configuration space  come with a price. Let us look at simple example.
The two-dimensional isotropic harmonic oscillator and two identical, spatially
distant, non-interacting one-dimensional oscillators have exactly the same
configuration space, $\mathbb{R}^2$, and, once the two one-dimensional
coordinates are laid side by side, exactly the same Hamiltonian. This
coincidence might seem to threaten the idea that configuration space carries
direct physical meaning: if the bare pair $(Q,H)$ cannot distinguish ``one
particle moving in a plane'' from ``two particles moving on separate, distant
lines,'' in what sense does configuration space represent physical positions at
all? 

This coincidence  persists at every
level of structure computable from $(Q,H)$ alone --- classical trajectories,
the symplectic structure, the quantum spectrum and its degeneracies.
Additional data is needed to tell the two systems
apart: 
\begin{enumerate}
   \item an embedding of $Q$ into physical space together with the action of
the spatial isometry group, which is geometrically realized in one case and
merely accidental in the other;
   \item statistics and particle-content data,
which matter only when exchange effects are actually observable;
   \item the
algebra of interactions that locality in physical space picks out as natural,
which governs how each system couples to a generic environment.
\end{enumerate}
Configuration space, taken on its own, is best understood as a
representation-dependent bookkeeping device rather than a self-standing arena
of physical positions: its physical reading requires further data as an
independent postulate, not something read off the space itself.

Nevertheless, in most cases, especially when $K$ (holonomic) constraints 
 exist it is very convenient to forget
of physical meaning of coordinates and consider configuration space as $ n=3 k -K $
dimensional manifold  $ Q $ with arbitrary coordinates $q_1,q_2,\ldots q_n$ on it.

Velocities of the system live in the tangent bundle $TQ$; the Lagrangian is
intrinsically a function $L : TQ\times\mathbb R \to \mathbb R$, stated without
reference to any coordinates at all. The Euler--Lagrange
equations~\eqref{lagr-eq} themselves have a coordinate-free meaning: they are
the flow of a specific vector field on $TQ$, picked out by the
Poincar\'e--Cartan structure that $L$ induces on $TQ$, satisfying the
second-order condition that this flow really does generate velocities
consistent with positions. What one writes down in coordinates
as~\eqref{lagr-eq} is the local expression of this single intrinsic object,
valid in every chart, looking different in each.

Passing (via the Legendre transform, where it is well defined) to
the momenta 
\begin{equation}
   \label {gen-mom}
   p_i = \partial L/\partial \dot q^i
\end{equation}
carries the
description to the cotangent bundle $T^*Q$, the phase space. This
space carries a canonical structure that exists with reference to no
coordinates whatsoever: the \emph{tautological} (or Liouville) 1-form
$\theta$, defined at a point $(q,p)\in T^*Q$ by $\theta(v) =
p\bigl(\pi_*v\bigr)$ for any tangent vector $v$ there, where $\pi:
T^*Q\to Q$ is the projection. Its exterior derivative $\omega =
d\theta$ is the canonical symplectic form on $T^*Q$: closed,
nondegenerate, and built entirely out of the bundle structure of
$T^*Q$ over $Q$, with no reference to a coordinate system anywhere in
its definition.

Hamilton's equations are then the local expression of a single
intrinsic equation,
\begin{equation}
  \iota_{X_H}\,\omega \;=\; dH,
  \label{hamiltonian-vf}
\end{equation}
which defines the Hamiltonian vector field $X_H$ for any smooth
function $H$ on phase space. In local coordinates $(q^i,p_i)$ in
which $\theta = p_i\,dq^i$ and $\omega$ takes its standard constant
form, \eqref{hamiltonian-vf} reproduces exactly the familiar
\begin{equation}
  \dot q^i \;=\; \frac{\partial H}{\partial p_i},
  \qquad
  \dot p_i \;=\; -\,\frac{\partial H}{\partial q^i}.
  \label{hamilton-coord}
\end{equation}
Coordinates in which $\omega$ takes this simple constant form are
called \emph{canonical}, or \emph{Darboux}, coordinates, after Gaston
Darboux, who proved in 1882 that such coordinates always exist, at
least locally, around every point of every symplectic
manifold. This is a striking fact with no analogue
in Riemannian geometry, where curvature is a genuine local invariant
obstructing the metric from being brought to a standard flat form
even at a single point's neighborhood. Symplectic manifolds have no
such local invariant: every symplectic manifold of a given dimension
looks, locally, exactly like every other one. This is precisely why
Hamilton's equations can always be brought to the same simple
form~\eqref{hamilton-coord}, no matter how complicated the
physical system --- the complexity of a hard mechanics problem lives
entirely in the choice of \emph{which} canonical coordinates to use,
never in whether they exist.


\paragraph{"Spin" in constant magnetic field}
\label{par:cm-spin}
We used to natural systems where hamiltonian is quadratic form over momenta and
phase space is $T^*Q$ 
but there are physical systems for which both statements are false.
Let us consider 
dynamic of point magnetic moment in constant magnetic field has not only
clear physical meaning but also significant practical value. I investigated
this problem (actually slightly more general one) in connection with control of MRI
scanner \cite{RPZ4},\cite{RPZ3},\cite{RPZ2},\cite{RPZ1}.


Hamiltonian is $ H=\bm M \cdot \bm B $ where $\bm M$ is magnetic moment 
and $ \bm B $ is constant magnetic field.
It is convenient to write $ \bm M=\frac{e}{m c}\bm S $ where $ \bm S $
is angular moment ("spin")  and  include gyromagnetic ratio into magnetic field
$ \bm \Omega = \frac{e}{m c}\bm B$. Components of $ \bm S $ are not
canonical variables, their Poisson brackets are
\begin{equation}
   \label {pb}
    \{S_i,S_j\}= \epsilon_{ijk} S_k
\end{equation}
However they define phase space (sphere)  and equation of motion:
\begin{equation}
   \label {sp1}
    \dot S_i=\{S_i,H\}
\end{equation} 
or 
\begin{equation}
   \label {sp2}
    \dot {\bm S} =-\bm S \times \bm \Omega
\end{equation}

In canonical variables 
\begin{equation}
   \label {sp3}
   \begin{split}
       S_x & = \sqrt{S^2-p^2} \cos q \\
       S_y & = \sqrt{S^2-p^2} \sin q \\
       S_z & = p,\quad -S\le p \le S
   \end{split}
\end{equation}
Let $ \bm \Omega =(0,0, \Omega) $ then $ H=p \Omega $ and canonical equations are:
\begin{equation}
   \label {sp4}
   \begin{split}
       \dot p & =0\\
       \dot q & = \Omega
   \end{split}
\end{equation}
Lagrangian $ L=0 $, so we can not get equations of motion as second order equation.

Moment $ p $ is integral of motion which is fine but velocity is not connected to 
moment at all, it is constant. For pair of points in phase space generally there is
no trajectory passable for given time. Variation problem simply can not be posed.

Hamiltonian $ H=p \Omega $ looks like that for harmonic oscillator in
action-angle variables and so we can  transform it to usual quadratic form, but
the condition $  -S\le p \le S$ makes this transformation pure formal.

\section{Constraints}
\label{sec:cm-const}
Constraints---restrictions on the positions and velocities a mechanical system may
assume---are present in practically every realistic model.   A classification may be established along several independent axes: the order of the
derivatives involved (geometric, kinematic, dynamic), integrability (holonomic,
semi-holonomic, nonholonomic, with the finer distinction between partially and completely
nonholonomic systems), explicit time dependence (scleronomic, rheonomic), sidedness
(bilateral, unilateral), the work done by the constraint reactions (ideal, non-ideal),
regularity, and, in the Hamiltonian setting, Dirac's primary/secondary and
first-class/second-class constraints.
Only ideal holonomic  constraints will be discussed here.

Let $q=(q^1,\dots,q^n)$ be generalized coordinates.
A system of constraints is \emph{holonomic} if it is
    equivalent to a system of finite
relations (geometric constraints)
\begin{equation}
   \label {holonomic-const}
   f_\beta(t,q)=0 \quad \beta =1,\dots m .
\end{equation}
 Note that since
 \[\df{f_\beta}{t}=\pd{f_\beta}{t}+\sum_{i=k}^n\pd{f_\beta}{q_k} \dot q_k=0 \]
every \emph{holonomic} constraint is can be seen as \emph{integrable} 
Pfaffian (linear over velocities) constraint:
\begin{equation}
  \sum_{i=1}^n a_i(q,t) \dot q^i + a_t(q,t) = 0,
  \label{pfaffian-const}
\end{equation}

\paragraph{d'Alembert--Lagrange principle and Lagrange multipliers}
Let $L(q,\dot q,t)=T-V$ and let $Q_i$ denote applied generalised forces not derived from
$V$. The \emph{d'Alembert--Lagrange principle} states that for \emph{ideal} constraints the
true motion satisfies
\begin{equation}\label{dalembert}
  \sum_{i=1}^n\Bigl(\frac{d}{dt}\pd L{\dot q^i}-\pd L{q^i}-Q_i\Bigr)\delta q^i=0
  \qquad\forall\,\delta q\in\mathcal V_{(t,q)}.
\end{equation}
Because $\mathcal V_{(t,q)}=\{\delta q:\sum a_{\alpha i}\delta q^i=0\}$, the bracket must lie in the
span of the constraint gradients, which gives the \emph{Lagrange equations of the first kind}
\begin{equation}
    \label{lagrange-first}
  \frac{d}{dt}\pd L{\dot q^i}-\pd L{q^i}=Q_i+\sum_{\alpha=1}^m\lambda_\alpha\,a_{\alpha i}(q,t),
  \qquad i=1,\dots,n,
\end{equation}
together with the $m$ constraint equations \eqref{pfaffian-const}. The $n+m$ equations determine
the $n$ functions $q^i(t)$ and the $m$ multipliers $\lambda_\alpha(t)$. The constraint force
is $R_i=\sum_\alpha\lambda_\alpha a_{\alpha i}$; in particular the multipliers carry
direct physical meaning (tensions, normal forces). For a holonomic constraint $f_\beta(q,t)=0$
one has $a_{\beta i}=\partial f_\beta/\partial q^i$.

\paragraph{Holonomic reduction.} If all constraints are holonomic one can eliminate
the multipliers \emph{a priori}: choose $n-m$ independent coordinates $s=(s^1,\dots,s^{n-m})$
with $q=q(s,t)$ solving $f_\beta=0$, substitute into $L$ to obtain $\tilde L(s,\dot s,t)$, and
write the ordinary Lagrange equations
\begin{equation*}
   \label {}
    \frac{d}{dt}\pd{\tilde L}{\dot s^r}-\pd{\tilde L}{s^r}=\tilde Q_r
\end{equation*}
 with no reactions at all.
This is the standard route and the main practical benefit of holonomic constraints.

Physically constraints are  the idealization build from the 
family of smooth models. A holonomic constraint $f=0$ is the limit
$\kappa\to\infty$ of a stiff potential $V_\kappa=\tfrac12\kappa f^2$. The limit
is delicate (see chapter \ref{lm}).

If the initial position lies on the surface and the initial velocity is tangent
to it, the motion converges to the d'Alembert motion on the surface. If the
initial velocity has a normal component, the normal motion is a fast
oscillation with a conserved action, and the limiting motion on the surface is
governed by the original Hamiltonian \emph{plus an extra potential} that
depends on that action. Rubin and Ungar \cite{RubinUngar1957}proved this for
codimension one; Takens \cite{Takens1980} treated the general case, under a
non-resonance condition on the frequencies of the normal oscillations. When
that condition fails (frequencies crossing along the surface), there may be no
differential equation on the surface governing the limiting motion at all.
Similarly, a nonholonomic constraint (not discussed here)  is the limit of a
large viscous friction $-\mu\,(a_i\dot q^i)\,a^i$ as $\mu\to\infty$. 

This is a clean illustration of the thesis that a theory is not a limit of
another but a chart on it: the same rod, with the same symbol $f=0$, belongs
to different physical charts depending on whether it is the limit of a stiff
spring or the limit of a viscous friction.

\section[Geometry vs Coordinates]{Geometric Elegance versus Coordinate Practice }
\label{sec:cm-geom-vs-coord}
The geometric form tells, in advance and for
free, which features of a set of equations of motion are physical
and which are artifacts of a chart. It explains why canonical
transformations exist and behave the way they do (Darboux's
theorem). It makes the relationship between symmetry and conservation
essentially automatic: a continuous symmetry is the flow of a vector
field preserving both $\omega$ and $H$, and when that symmetry comes
from the action of a Lie group, the associated conserved quantity is
the group's \emph{momentum map} --- the modern, coordinate-free
statement of Noether's theorem.
It also exposes genuinely global, topological facts that no single
coordinate patch can see: the phase space of a rigid body, for
instance, is the cotangent bundle of the rotation group $SO(3)$,
which cannot be covered by a single nonsingular coordinate chart at
all --- the root cause of gimbal lock, which textbooks often present
as an engineering nuisance but which is, geometrically, an
unavoidable topological obstruction. A coordinate singularity, on
this view, is never a physical singularity: it is only ever the
failure of one particular chart, repaired by changing coordinates or,
better, by working with the coordinate-free structure directly, where
the question of a "singularity" never even arises.

The coordinate form is actually the only form in which a
problem is actually solved, in closed form or on a computer, and the
only form that connects directly to a measurement, which is always
reported as a number relative to some frame. Darboux's theorem
guarantees good coordinates exist near any point; it says nothing
whatsoever about how to find the ones adapted to a \emph{specific}
Hamiltonian's structure --- action-angle variables for an integrable
system, say, which is frequently the entire technical content of
actually solving a hard problem in mechanics. No amount of abstract
symplectic geometry does that work for you.

Seen from the perspective of the wider project of this book, the two
forms sit on opposite sides of its recurring divide: a coordinate
expression is what gets compared, number for number, against
measurement, and so belongs to the "real world" side of the ledger;
the manifold, the connection, the symplectic form are pieces of
"ideal world" mathematics, whose entire justification is that they
correctly organize and predict the coordinate-level facts, and
nothing more. Neither side is dispensable, and neither is more
fundamental than the other in any sense that would let one be
discarded.

A single example run in both languages makes the trade-off concrete.
A pendulum of length $\ell$ and bob mass $m$, free to swing in any
direction, has configuration space the sphere of radius $\ell$
centered on the pivot.

\emph{Coordinate form.} Using the polar angle $\theta$ from the
downward vertical and the azimuth $\phi$, the Lagrangian is
\begin{equation}
  L \;=\; \tfrac12 m\ell^2\bigl(\dot\theta^2 + \sin^2\theta\,
  \dot\phi^2\bigr) \;+\; mg\ell\cos\theta .
  \label{sph-pendulum-L}
\end{equation}
The coordinate $\phi$ does not appear in $L$ itself, so its conjugate
momentum
\begin{equation}
  p_\phi \;=\; m\ell^2\sin^2\theta\,\dot\phi
  \label{sph-pendulum-pphi}
\end{equation}
is conserved --- it is the angular momentum about the vertical axis
--- and the problem reduces to an effective one-dimensional motion in
$\theta$ alone, with effective potential $V_{\mathrm{eff}}(\theta) =
mg\ell\cos\theta + p_\phi^2/(2m\ell^2\sin^2\theta)$. This is the form
in which the problem is actually solved, e.g.\ by quadrature or
numerically.

\emph{Geometric form.} Configuration space is the sphere $S^2$
itself, with the circle group $SO(2)$ acting on it by rotation about
the vertical axis, an action that preserves the Lagrangian. The
conservation of $p_\phi$ is nothing but the momentum map for this
$SO(2)$ symmetry, an instance of Noether's theorem requiring no
mention of $\theta,\phi$ at all. And the coordinate singularities of
$(\theta,\phi)$ at the poles $\theta=0,\pi$ --- where $\phi$ becomes
undefined and equation~\eqref{sph-pendulum-pphi} degenerates ---
are transparently artifacts of this particular chart: $S^2$ has no
distinguished points and no singularity there whatsoever. A pendulum
hanging straight down, or balanced straight up, is a perfectly smooth
physical configuration; it is only the coordinate system built from
$(\theta,\phi)$ that breaks down at exactly those two points, for the
same reason that all spherical coordinate systems break down at their
poles.

The two languages are not competitors; they specialize in different
questions. The geometric form is the place to ask what is really
happening, structurally, and what is guaranteed to remain true no
matter which coordinates one happens to choose. The coordinate form
is the place --- indeed the only place --- to ask what number a
system actually produces. A physicist who works only in coordinates
is at permanent risk of mistaking an artifact of a chosen chart for a
physical effect, exactly as happened, for a good part of the
twentieth century, with the coordinate singularity of the
Schwarzschild solution at its horizon: a singularity of the
particular $(t,r,\theta,\phi)$ coordinates originally used to write
the metric down, mistaken for decades for a genuine physical
singularity of spacetime itself, and only cleared up by passing to
coordinates (Eddington--Finkelstein, Kruskal) in which the horizon is
manifestly perfectly smooth. A physicist who
works only in the geometric language, on the other hand, has an
elegant theory of everything and the means to compute nothing. The
working practice of mechanics has always required both, spoken
fluently, and switched between as the problem at hand demands.

